%% sections/03-methodology.tex

\section{Methodology: The Amendment Pipeline}
\label{sec:methodology}

Innovation-cluster-driven rubric amendment is a three-stage process that operates
continuously throughout a corpus of evaluated sessions. It requires no separate
review cycle: the pipeline runs as part of the evaluation process itself. This
section describes each stage in full operational detail, including the
decision criteria, data formats, and constraints that make the mechanism
replicable.

\subsection{Stage 1: Innovation Logging}

An innovation log entry is created immediately after the session-gate scoring
stage when the gate reviewer identifies an outcome that satisfies three criteria:

\begin{itemize}
  \item the outcome is clearly a quality achievement, not an evaluative failure
        or an artifact of the session's particular content;
  \item the outcome is not described by any existing rubric anchor at the
        applicable version (i.e., a scorer consulting only the rubric would have
        no specific guidance for this pattern);
  \item the outcome could plausibly recur in future sessions under some
        recognizable condition (not purely an idiosyncrasy of a single session).
\end{itemize}

Each entry has five fields:

\begin{description}
  \item[Technique:] A short noun-phrase naming the observed pattern
    (e.g., ``Cross-character mechanical compound,'' ``Sustained behavioral
    pattern as choice,'' ``Portent-as-directed-foreknowledge'').

  \item[Dimension:] The primary rubric dimension the innovation is most
    relevant to. Cross-dimension relevance is noted as secondary but does not
    affect the clustering trigger.

  \item[Evidence:] A direct quotation from the session log with scene
    reference, plus a brief analytical note explaining why the existing anchors
    did not capture this outcome.

  \item[Scope:] One of three values --- \textit{party-specific} (likely to
    recur only with this party's composition or grief-paragraph configuration),
    \textit{adventure-specific} (likely to recur with this adventure design
    pattern), or \textit{universal} (likely to recur with any party or
    adventure that produces the relevant conditions).

  \item[Status:] One of: \texttt{logged} $\to$ \texttt{proposed} $\to$
    \texttt{adopted (vX.Y)} or \texttt{reviewed-not-adopted}.
\end{description}

The log is append-only. No entry is deleted or modified after creation; status
updates are the only permitted modification. This preserves the full history of
what the rubric did not capture and when.

\subsection{Stage 2: Clustering and Threshold Check}

After each session, the full set of entries with status \texttt{logged} is
scanned for dimension clusters. The trigger condition is:

\[
  \bigl|\{e \in \mathcal{L} : e.\mathrm{dim} = d
          \;\wedge\; e.\mathrm{status} = \texttt{logged}\}\bigr| \;\geq\; 2
\]

for any dimension $d$, where $\mathcal{L}$ is the current log. When the
threshold fires, an amendment proposal is drafted immediately.

Clustering does not require entries to be semantically identical: it requires
that they tag the same dimension from different session instances. In practice,
dimension co-tags are a reliable proxy for thematic coherence. Each dimension
names a conceptually bounded scoring target (e.g., Atmospheric Landing scores
how emotional and atmospheric intent is realized in play); entries that share a
dimension tag consistently address the same underlying quality phenomenon from
different session contexts.

The 2-entry threshold is intentionally conservative. It rules out single-instance
patterns --- which may be session-specific artifacts --- while permitting rapid
response to genuinely recurring phenomena. In the Marathon corpus, all 15
amendment-triggering clusters met the 2-entry threshold; no cluster required
a higher threshold to achieve thematic coherence (though two clusters of 3 and
one cluster of 4 occurred and reinforced the resulting amendments).

\paragraph{Formal threshold definition.}
Let $\mathcal{I}_d$ be the set of logged innovations with primary dimension tag $d$
and status \texttt{logged}. An amendment proposal for dimension $d$ is triggered when
$|\mathcal{I}_d| \geq 2$ AND the two constituent innovations share a thematic cluster
--- that is, they address the same underlying quality phenomenon from different session
instances, even if they describe the phenomenon in distinct surface forms. Dimension
co-tagging is the primary clustering criterion; the thematic-cluster requirement rules
out pairs of entries that share a dimension tag but describe unrelated phenomena within
that dimension's scope.

\paragraph{False-positive rate in the corpus.}
Of the 13 ratified amendments in the Marathon corpus, zero were later reversed. No
amendment that passed the 2-instance-plus-thematic-cluster criterion was subsequently
found to describe a spurious or non-recurring pattern. The threshold, in combination
with the thematic-clustering requirement, produced no false positives across the full
corpus.

\paragraph{Alternative threshold considered.}
A 3-instance threshold was considered during the rubric's development. Under a
3-instance rule, amendment v1.1 (constituent innovations I-S01-04 and I-S01-07;
exactly 2 instances) and amendment v1.9 (constituent innovations I-C4-S1-01 and
I-C4-S1-02; exactly 2 innovations in the AL cluster) would both have been delayed
pending a third observation. Both proved valid: v1.1's cross-dimensional harmlessness
anchor and v1.9's grief-paragraph situational resonance anchor each produced no
subsequent reversals and were evidenced in later sessions without amendment. The
2-instance threshold is aggressive relative to the 3-instance alternative, but the
corpus provides no evidence that the additional speed produced erroneous amendments.
A calibration study comparing 2-instance vs.\ 3-instance thresholds across held-out
sessions would formalize this parameter; we treat it as a design choice pending
validation.

\subsection{Stage 3: Amendment Adoption}

An amendment proposal documents: (a) the constituent innovation entries,
(b) the dimension being amended, (c) the proposed anchor text, and (d) a
rationale explaining why the existing anchors did not cover the observed pattern.
Proposals are reviewed against two criteria:

\begin{enumerate}
  \item The described pattern is \textit{genuinely novel}: it is not captured
        by any existing anchor even under a charitable reading.
  \item The proposed anchor text is \textit{sufficiently specific}: a future
        evaluator consulting only the rubric would assign consistent scores on
        sessions where the pattern appears.
\end{enumerate}

Proposals that pass both criteria are adopted: the rubric version is incremented,
constituent innovations are marked \texttt{adopted (vX.Y)}, and the new anchor
is added to the rubric document. The version increment follows a minor-version
convention: v1.X increments to v1.(X+1). No major-version bumps occurred in the
Marathon corpus because all amendments were additive (new anchors added) rather
than corrective (existing anchors modified).

In the Marathon corpus, all 15 proposals were adopted. No proposals were marked
reviewed-not-adopted. This likely reflects the conservative drafting standard:
proposals were only drafted when the cluster was unambiguous and the textual
evidence was strong. It is reasonable to expect that a more liberal logging
standard would produce some reviewed-not-adopted outcomes.

\subsection{The Forward-Only Constraint}

All amendments apply forward-only: sessions scored before an amendment cannot
be re-scored under the new anchor. This creates version-windows in which all
scores are comparable: sessions scored under v1.3 are directly comparable to
each other, and cross-version comparisons require version as an explicit covariate.

The forward-only constraint serves two functions. First, it prevents retroactive
inflation: if amendments could be applied retroactively, every improvement in
rubric expressiveness would immediately improve all prior scores, conflating
measurement improvement with quality improvement. Second, it preserves the
operational record: the innovation log accurately reflects what the rubric
\textit{did not capture} at the time of each session, which is the primary
empirical record of the mechanism.

\subsection{Complete Amendment History}

Table~\ref{tab:full-history} provides the complete amendment history across all
13 rubric versions, including the constituent innovations and the dimension
affected at each step.

\begin{table*}[t]
\caption{Complete Rubric Amendment History: v1.0 Through v1.12}
\label{tab:full-history}
\setlength{\tabcolsep}{5pt}
\begin{tabular}{@{}p{0.9cm}p{1.2cm}p{3.0cm}p{3.2cm}p{5.0cm}@{}}
\toprule
Version & Dim. & Constituent Innovations & Session(s) & Amendment Description \\
\midrule
v1.0 & --- & --- & Pre-S01 & Initial rubric. Eight dimensions:
  Engagement (EG), Mechanical Fairness (MF), Pacing (PA), Character Agency (CA),
  Module Fidelity (MFid), Atmospheric Landing (AL), Surprise (SU),
  Table Readiness (TR). \\
\addlinespace[2pt]
v1.1 & MF & I-S01-04, I-S01-07 & S01 & Cross-dimensional harmlessness anchor.
  A mechanical interaction that achieves high MF in its primary dimension may
  not depress scores in an unrelated dimension; cross-dimensional neutrality is
  the baseline. \\
\addlinespace[2pt]
v1.1 & AL & I-S01-02, I-S01-10, I-S01-12 & S01 & Per-PC reception tracking
  (two-column finder/receiver structure); inter-PC chain-reaction scores at the
  10-band when an atmospheric beat received by one PC visibly propagates to a
  second PC. \\
\addlinespace[2pt]
v1.2 & MFid & I-S01-05, I-S01-06, I-S01-08, I-S01-11 & S01--S02 & Anchor-level
  single-gate failure. Module Fidelity can score 8 at anchor level while failing
  to deliver a single designed symptom; the 10-band requires all designed anchors
  present and legible, not merely present. \\
\addlinespace[2pt]
v1.2 & AL & I-S02-01, I-S02-03 & S02 & Finder/receiver divergence: a PC who
  discovers an atmospheric element may not be the PC who receives it emotionally;
  both roles are tracked. Silent reception (no verbal acknowledgment) now
  explicitly scores at the 9-band when the log provides behavioral evidence. \\
\addlinespace[2pt]
v1.3 & MFid & I-S03-01, I-S03-02, I-S03-04 & S03 & Positional-geometry failure.
  A combat encounter that fails to instantiate a designed terrain-challenge is
  a Module Fidelity miss even if tactical gameplay was otherwise high quality.
  Anchor-level symptoms must include spatial/positional elements when the module
  specifies them. \\
\addlinespace[2pt]
v1.4 & MF & I-S03-03, I-S04-01 & S03--S04 & Cross-character mechanical compound.
  A cross-PC mechanical interaction that enhances a second PC's roll (e.g.,
  Bless cast to support a Religion check) scores at the 10-band if voiced in
  character. Portent-as-benediction is the canonical exemplar. \\
\addlinespace[2pt]
v1.4 & AL & I-S04-02, I-S05-01 & S04--S05 & Multi-scene silent-reception chain.
  Reception that persists across two or more scenes without verbal acknowledgment,
  evidenced by behavioral callbacks (a carried object referenced, an earlier
  statement echoed), scores at the 10-band for the chain, not only the originating
  scene. \\
\addlinespace[2pt]
v1.5 & AL & I-S07-01, I-C2-S01-01, I-C2-S02-03 & S07, C2-S01--S02 &
  Cross-session act-without-announcement. A PC behavioral pattern that expresses
  a received atmospheric element across sessions --- rather than within a single
  session --- scores as a campaign-level AL achievement at the 10-band. The
  session log must cite the originating session. \\
\addlinespace[2pt]
v1.6 & AL & I-S07-02, I-C2-S04-01 & S07, C2-S04 & Simultaneous arc-convergence.
  When two or more PCs reach an emotional arc-completion point at the same scene,
  the convergence itself scores as an AL event independent of each individual
  arc-completion; the evaluator notes the convergence explicitly. \\
\addlinespace[2pt]
v1.7 & AL & I-C3-S01-01, I-C3-S01-02 & C3-S01 & Framework-failure arc-activation.
  A PC whose arc-completion is triggered by the failure of an in-world authority
  figure or institutional structure (rather than by discovery or narrative gift)
  scores at the 10-band for activation mode; the framework failure is the
  atmospheric delivery mechanism. \\
\addlinespace[2pt]
v1.8 & AL & I-C3-S03-01, I-C3-S04-01 & C3-S03--S04 & Arc-activation through
  behavioral stance. A PC who occupies a particular behavioral stance (restraint,
  witness, deliberate inaction) across a scene or encounter and whose arc advances
  through that stance --- rather than through explicit speech or narrative gift ---
  scores at the 10-band for the activation. \\
\addlinespace[2pt]
v1.9 & AL, CA & I-C4-S1-01, I-C4-S1-02, I-C4-S1-03, I-C4-S1-04 & C4-S01 &
  Dual amendment. (AL) Grief-paragraph situational resonance: when a PC's
  grief-paragraph commitment directly shapes their resource decisions in a session
  with a sustained-pressure mechanic, the resonance scores as AL achievement at
  the 9-band. (CA) Sustained behavioral pattern as choice: a PC maintaining a
  consistent behavioral pattern across 3+ scenes --- rather than making a single
  visible choice --- now qualifies as a CA achievement at the 9-band. \\
\addlinespace[2pt]
v1.10 & CA & I-C4-S2-01, I-C4-S3-01 & C4-S02--S03 & Character-register
  conversion. A PC who translates between two character registers within a single
  scene (e.g., moving from tactical-practical to grief-honest speech without
  narrative prompt) scores at the 10-band for CA; the conversion must be
  unprompted and evidenced by direct quotation. \\
\addlinespace[2pt]
v1.11 & MF & I-C4-S3-02, I-C4-S4-01 & C4-S03--S04 & Portent-as-directed-fore\-knowledge.
  A Portent die held and deployed to a named ally rather than retained for personal
  use --- with explicit in-voice framing of the delegation --- scores at the 10-band
  for MF. The delegation must be voiced; silent deployment does not qualify. \\
\addlinespace[2pt]
v1.12 & AL & I-C4-S2-02, I-C4-S3-03 & C4-S02--S03 & Intra-session restraint as
  reception. A PC who receives an atmospheric element and deliberately defers its
  expression to a later scene within the same session --- placing an object aside,
  declining to speak, returning to a threshold --- scores at the 9-band for the
  deferred expression rather than the originating beat. \\
\bottomrule
\end{tabular}
\end{table*}

The table reveals several structural features of the amendment history.
First, v1.1 was a compound amendment: it incremented the version once but
addressed two distinct dimension clusters (MF and AL) simultaneously. In
subsequent amendments, each version increment addressed one cluster, with
the exception of v1.9 (a dual amendment addressing an AL cluster and a
CA cluster that fired simultaneously from the same session). Second, the
gap between v1.8 and v1.9 --- both running through C3's 7 sessions --- shows
the rubric's stability once it has absorbed a campaign archetype. Third,
v1.12 has been stable across all of C5, C6, and C7, which together span
21 sessions (7 per campaign).

\subsection{Operationalizing the Pipeline in a New Evaluation Context}

To implement the mechanism in a different creative evaluation context, a
practitioner requires:

\begin{enumerate}
  \item A \textbf{session log format} that records specific outcomes with
        sufficient textual granularity to support evidence citation. Vague
        summaries are insufficient; the innovation log requires direct quotation
        or paraphrase with scene reference.
  \item A \textbf{rubric with named dimensions} and explicit anchor descriptions
        at the score bands being used. The clustering trigger operates on
        dimension tags; without named dimensions, clustering cannot occur.
  \item A \textbf{forward-only versioning convention}, enforced consistently
        from the first session. Retroactive scoring creates version-window
        ambiguity that undermines the mechanism's primary statistical property.
\end{enumerate}

The innovation log schema is available in the Marathon workshop repository
alongside the rubric format and amendment proposal template.
